% SGA 3，报告 VII_B，英文重建本——第 1 轮（中文译本）。
% 范围：依据原文件第 6--8 页，从 0.2.B 证明后至 0.2.2 结束；止于 0.3 之前。
% 合并阅读版页码： 519--521. 再版印刷页码： 508--510.
% 本范围内的原版页边定位： 481--482.
% 依据： Polo--Gille Exp7B-13oct24.pdf,
% SHA-256 D303690E0195654BBE7C3DAA5D89EA23E24CD709A3575217234AC19AE8836694.
% 曾参阅 Jacob Reinhold 的 CC BY 4.0 Markdown，作为对比版本来源。

\medskip
设 \((\mathrm{Ab})\)\footnote{\label{sga3:VIIB:nde:10}编者在本段中插入了\SGARefLink{sga3:VIIB:0.2.E}{命题 0.2.E} 和
\SGARefLink{sga3:VIIB:0.2.F}{推论 0.2.F}；它们将在
\SGARefLink{sga3:VIIB:0.2.2}{0.2.2} 中用到。在原版中，这些结果位于
\SGARefLink{sga3:VIIB:0.3}{0.3}。} 为阿贝尔群范畴，设 \(\operatorname{LF}(A)\) 为
\(\operatorname{PC}(A)\) 中由有限长度的对象构成的满子范畴。对
\(\operatorname{PC}(A)\) 的每个对象 \(M\)，以 \(h_c^M\) 表示函子
\[
  \operatorname{LF}(A)\longrightarrow(\mathrm{Ab}),
  \qquad
  N\longmapsto\operatorname{Hom}_c(M,N).
\]
由 \textup{[CA]}，\S~II.4，定理~1、引理~4 和推论~1，得到下列结果。\footnote{\label{sga3:VIIB:nde:11}事实上，\(\operatorname{PC}(A)\) 有一个 Artin 余生成元集合，其中的滤过投射极限是正合的，
而 \(\operatorname{LF}(A)\) 是 Artin 对象子范畴。因此，对偶范畴
\(\operatorname{PC}(A)^\circ\) 有一个 Noether 生成元集合，且其中的滤过归纳极限是正合的。
由 \textup{[CA]}，\S~II.4，定理~1 的证明，函子
\(M\mapsto\operatorname{Hom}_c(M,-)\) 是从 \(\operatorname{PC}(A)\) 到
\(\operatorname{Lex}(\operatorname{LF}(A),(\mathrm{Ab}))\) 的反等价。
同样，\textit{同处}的引理~4 和推论~1 陈述了与
\SGARefLink{sga3:VIIB:0.2.F}{推论~0.2.F} 中的结果“对偶”的结果。
“直接”的证明见 \textup{[DG70]}，\S~V.2，定理~3.1、引理~3.5 以及推论~3.3 和~3.4。}

\medskip
\noindent\textbf{命题 0.2.E。}
\phantomsection
\label{sga3:VIIB:0.2.E}
\textit{函子 \(M\mapsto h_c^M\) 给出从 \(\operatorname{PC}(A)\) 到左正合函子
\(\operatorname{LF}(A)\longrightarrow(\mathrm{Ab})\) 的范畴}
\[
  \operatorname{Lex}\bigl(\operatorname{LF}(A),(\mathrm{Ab})\bigr)
\]
\textit{的反等价。}

\medskip
\noindent\textbf{推论 0.2.F。}
\phantomsection
\label{sga3:VIIB:0.2.F}
{\itshape
\begin{enumerate}
  \item[\textup{(i)}]
  \(\operatorname{PC}(A)\) 的对象 \(P\) 是投射对象，当且仅当函子
  \(h_c^P\) 正合；即当且仅当函子 \(\operatorname{Hom}_c(P,-)\) 在
  \(\operatorname{LF}(A)\) 上正合。

  \item[\textup{\SGAInlineRefTarget{sga3:VIIB:0.2.F:item:ii}(ii)}]
  设 \((M_i)\) 是 \(\operatorname{PC}(A)\) 中对象的滤过投射系\footnote{\label{sga3:VIIB:nde:12}由于每个无穷积都是有限积的滤过投射极限，
  每个与滤过投射极限交换的加性函子也都与无穷积交换。} 。对 \(\operatorname{LF}(A)\) 的每个对象 \(N\)，有关于 \(N\) 函子性的同构
  \[
    \operatorname{Hom}_c\!\left(\varprojlim_i M_i,N\right)
    \simeq
    \varinjlim_i\operatorname{Hom}_c(M_i,N).
  \]

  \item[\textup{\SGAInlineRefTarget{sga3:VIIB:0.2.F:item:iii}(iii)}]
  \(\operatorname{PC}(A)\) 中投射对象的每个滤过投射极限和每个积\footnotemark[\value{footnote}] 都是 \(\operatorname{PC}(A)\) 的投射对象。
\end{enumerate}}

最后，由\SGARefLink{sga3:VIIB:0.2.F}{推论~0.2.F} 得到下列结果。

\medskip
\noindent\textbf{推论 0.2.G。}
\phantomsection
\label{sga3:VIIB:0.2.G}
\textit{设 \((M_i)_{i\in I}\) 是 \(\operatorname{PC}(A)\) 的一族对象。
则 \(\prod_{i\in I}M_i\) 是投射对象，当且仅当每个 \(M_i\) 都是投射对象。}

事实上，对每个 \(N\in\operatorname{Ob}\operatorname{LF}(A)\)，有
\[
  \operatorname{Hom}_c\!\left(\prod_i M_i,N\right)
  \simeq
  \bigoplus_i\operatorname{Hom}_c(M_i,N).
\]

\medskip
\noindent\textbf{0.2.1.}
\phantomsection
\label{sga3:VIIB:0.2.1}
\(A\) 的每个局部分量 \(A_{\mathfrak m}\) 都是 \(A\) 的直因子，
从而是 \(\operatorname{PC}(A)\) 的投射对象（\(A\) 显然是投射对象）。此外，
\[
  S_{\mathfrak m}
  =
  A_{\mathfrak m}/\mathfrak mA_{\mathfrak m}
\]
是 \(A_{\mathfrak m}\) 的唯一单商，所以 \(A_{\mathfrak m}\) 不可分解。
另一方面，\(\operatorname{PC}(A)\) 的每个单对象都同构于唯一的
\(S_{\mathfrak m}\)。由 \textup{[CA]}，IV，\S~3，定理~3 的推论~1，\footnote{\label{sga3:VIIB:nde:13}这里指的是\textit{同处} \S~2 的定理~2 及 \S~1 的命题~2 中建立的“对偶”陈述。
“直接”的证明见 \textup{[DG70]}，V，\S~2，定理~4.5 和例~4.6(b)。} 因此有：

\begin{samepage}
\medskip
\noindent\textbf{命题。}
\phantomsection
\label{sga3:VIIB:0.2.1:projectives}
{\itshape
\begin{enumerate}
  \item[\textup{(i)}]
  \(\operatorname{PC}(A)\) 的每个投射对象都是不可分解投射对象的直积，
  这些因子在同构意义下唯一确定。

  \item[\textup{\SGAInlineRefTarget{sga3:VIIB:0.2.1:item:ii}(ii)}]
  每个不可分解投射对象都同构于唯一的 \(A_{\mathfrak m}\)，其中
  \(\mathfrak m\in\Upsilon(A)\)。
\end{enumerate}}
\end{samepage}

\medskip
\noindent\textbf{定义。}
\phantomsection
\label{sga3:VIIB:0.2.1:topologically-free}
若伪紧 \(A\)-模 \(M\) 同构于一族 \(A\) 的副本 \((A_i)\) 的积，
则称 \(M\) 是\emph{拓扑自由}的。

在这种情形下，若将 \(A_i\) 的单位元映为 \(m_i\) 的 \(A\)-线性映射
\(A_i\longrightarrow M\) 延拓为同构
\[
  \prod_i A_i\xrightarrow{\ \sim\ }M.
\]
则称 \((m_i)\) 为 \(M\) 的一个\emph{伪基}。

\medskip
\noindent\textbf{0.2.2.}
\phantomsection
\label{sga3:VIIB:0.2.2}
\footnote{\label{sga3:VIIB:nde:14}编者详细展开了本段的结果，它们在下文起着重要作用；参见\
\SGARefLink{sga3:VIIB:1.2.3}{1.2.3}、\SGARefLink{sga3:VIIB:1.3.5}{1.3.5}、
\SGARefLink{sga3:VIIB:2.2.1}{2.2.1} 等。}
若 \(M\) 是伪紧 \(A\)-模，我们以 \(M^\dagger\) 表示不带拓扑的 \(A\)-模
\[
  \operatorname{Hom}_c(M,A).
\]

反之，若 \(N\) 是 \(A\)-模，则以
\[
  N^*=\operatorname{Hom}_A(N,A)
\]
表示其对偶，并赋予逐点收敛拓扑。因此，\(N^*\) 中 \(0\) 的一个邻域基由下列子模构成
\[
  \mathcal V(x,\mathfrak l)
  =
  \{f\in N^*\mid f(x)\in\mathfrak l\},
\]
其中 \(x\in N\)，且 \(\mathfrak l\) 是 \(A\) 的开理想。

这使 \(N^*\) 成为伪紧 \(A\)-模。事实上，先设 \(N=A\)。这时
\(N^*=A\)，并赋予其作为伪紧环的拓扑。若 \(N\) 是自由模 \(A^{(I)}\)，
则 \(N^*\) 是赋予积拓扑的积 \(A^I\)。此外，对于每个态射
\(\varphi:N_1\longrightarrow N_2\)，其转置
\[
  {}^t\!\varphi:N_2^*\longrightarrow N_1^*
\]
是连续的，因为 \(N_1^*\) 的子模 \(\mathcal V(x_1,\mathfrak l)\)
在 \({}^t\!\varphi\) 下的逆像恰是 \(N_2^*\) 的子模
\(\mathcal V(\varphi(x_1),\mathfrak l)\)。对任意 \(N\)，选取一个呈示
\[
  A^{(J)}
  \xrightarrow{\ \varphi\ }
  A^{(I)}
  \xrightarrow{\ \pi\ }
  N
  \longrightarrow
  0.
\]
则 \(N^*\) 是连续态射
\[
  {}^t\!\varphi:A^I\longrightarrow A^J,
\]
的核，因此 \(N^*\) 是伪紧的。

当 \(A\) 是 Artin 环时（这时上面可以取 \(\mathfrak l=0\)），
\SGARefLink{sga3:VIIB:0.2.F}{推论~0.2.F} 给出：

\medskip
\noindent\textbf{命题。}
\phantomsection
\label{sga3:VIIB:0.2.2:artinian-duality}
\textit{当 \(A\) 是 Artin 环时，函子}
\[
  P\longmapsto P^\dagger
  \qquad\textit{和}\qquad
  Q\longmapsto Q^*,
\]
\textit{在投射伪紧 \(A\)-模范畴与投射 \(A\)-模范畴之间建立反等价，其中
\(P\) 是 \(\operatorname{PC}(A)\) 的投射对象，而 \(Q\) 是投射 \(A\)-模。}\footnote{\label{sga3:VIIB:nde:15}还应回顾，在 Artin 环上，一个模是投射模当且仅当它是平坦模；
例如见报告 VI\textsubscript{B}，编者注~(54)，或下文
\SGARefLink{sga3:VIIB:0.3.7}{0.3.7}。}

特别地，当 \(A\) 是域 \(k\) 时，函子
\[
  P\longmapsto P^\dagger
\]
是所有伪紧 \(k\)-模（也称为\emph{线性紧 \(k\)-向量空间}）
构成的范畴与 \(k\)-向量空间范畴之间的反等价。\footnote{
\label{sga3:VIIB:nde:16}一族线性紧 \(k\)-向量空间 \((V_i)_{i\in I}\) 在 \(\operatorname{PC}(k)\) 中的直和
是 \(\left(\prod_{i\in I}V_i^\dagger\right)^*\)。}

% 下一断点： Exp7B 原文件页码 8, 合并阅读版页码 521,
% 再版印刷页码 510, 紧接其前： 0.3.
